% HAND-WRITTEN fragment -- not generated by maestro2tex, but placed by it via the
% `order` list in configs/anatop_tb_controlamp_cell.json.
%
% Numbers here come from the generated tables that follow: tab_ctrlamp_cell_corners,
% tab_ctrlamp_cell_loop, and the Monte Carlo tables. Re-read them from those tables
% if the run is repeated.
%
% Source runs: Interactive.158 (corners) and Plan.4.Run.2 (Monte Carlo), both
% archived under simarchive/anatop_tb_controlamp/cm_2026-07-27/.

\subsection{Potentiostat Control Amplifier}\label{ss:controlamp}

The control amplifier is amplifier A1 of the potentiostat. It holds the electrochemical
cell at a commanded potential: it drives the counter electrode (CE) from its output,
takes its feedback from the reference electrode (RE) on the inverting input, and takes
its set point from the reference DAC on the non-inverting input. Whatever current the
cell needs flows through the amplifier's output, in whichever direction the measurement
demands, so the amplifier must source and sink equally well. One instance is used per
channel; all four share the bias generator of Section~\ref{ss:biasgen-local}.

The topology is a \emph{symmetric current-mirror OTA} with a PMOS input pair
(Figure~\ref{fig:ctrlamp-core}). The input pair splits the tail current between two unit
diode loads; each half is then restated by a mirror four times larger, one directly at
the output and one folded through a PMOS mirror at node~\texttt{B}. Three properties
follow, and they are the reasons for the choice:

\begin{itemize}
\item \textbf{Source and sink are equal by construction.} The output current in either
direction is the same mirror ratio applied to the same tail current, so the two
directions match to the accuracy of the mirrors, \SI{6.4}{\percent} in Monte Carlo,
against roughly \(4{:}1\) for the folded-cascode amplifier used in the previous
revision.
\item \textbf{Peak output current is set by design, not by headroom.} It is the mirror
ratio times the tail current, \(4 \times \SI{5.85}{\micro\ampere} =
\SI{23.4}{\micro\ampere}\), which matches the \SIrange{23}{24}{\micro\ampere} measured.
\item \textbf{Current efficiency is high.} Quiescent supply is \((1+K)\,I_{\mathrm{tail}}\)
for mirror ratio \(K\), so \SI{80}{\percent} of the current drawn is available at the
output. The amplifier draws \SI{29.2}{\micro\ampere}.
\end{itemize}

The amplifier has a single high-impedance output node and no internal compensation
capacitor; the load sets the dominant pole. The MIM capacitors across the mirrors damp
the mirror poles. A PMOS
input pair places the input common-mode range against \texttt{vss} rather than the
supply, which is what a potentiostat holding an electrode near ground requires. The
\texttt{En} input powers the amplifier down by switching its three bias rails.

The characterisation below is schematic-level, with no extracted parasitics. Two
testbenches are used and the distinction between them matters when reading the numbers:

\begin{itemize}
\item \texttt{CellDriveTB} loads the amplifier with the \texttt{randles\_cell} model of
Section~\ref{ss:analog-models} and closes the loop from RE, which is how the amplifier
is actually used. Every drive, compliance and accuracy number comes from here.
\item \texttt{UnityFeedbackWBiasTB} and \texttt{StepUnityFeedbackWBiasTB} load it with
\SI{10}{\pico\farad} only. These give the small-signal and large-signal figures of the
amplifier itself. \textbf{Their gain figures are not the gain available at the
electrode:} a resistive cell loads the high-impedance output node and reduces loop gain
from \SI{76}{\decibel} to about \SI{8}{\decibel}
(Tables~\ref{tab:ctrlamp-mc-stability} and~\ref{tab:ctrlamp-cell-loop}). Regulation
accuracy follows the latter.
\end{itemize}
